\documentclass[a4paper,12pt]{article}

\usepackage[ngerman]{babel}
\usepackage[T1]{fontenc}
\usepackage[utf8]{inputenc}
\usepackage[a4paper,margin=2.5cm]{geometry}
\usepackage{amsmath,amssymb}
\usepackage{tikz}
\usepackage{xcolor}
\usepackage{enumitem}
\usepackage{lmodern}

\setlength{\parindent}{0pt}
\setlength{\parskip}{0.5em}

\begin{document}

{\LARGE\textcolor{red}{Einführung in die Theoretische Physik I}}

\vspace{1em}

{\Large\textcolor{red}{1. Vektoren}}

{\Large\textcolor{red}{1.1 Vektorrechnung}}

Betrachte Parallelverschiebungen (Translationen) der Punkte des 3D Anschaungsraumes,

\[
P_1=(x_1,y_1,z_1) \longrightarrow P_2=(x_2,y_2,z_2)
\]

Eigenschaften und Rechenoperationen anschaulich klar.

\[
\to \text{ Verallgemeinerung auf abstrakten Vektorbegriff im n-dimensionalen Raum.}
\]

\begin{center}
\begin{tikzpicture}[scale=1.2]
    \fill[red] (0,0) circle (1.5pt);
    \node[below left, red] at (0,0) {$P_1$};

    \fill[red] (3.5,1.2) circle (1.5pt);
    \node[right, red] at (3.5,1.2) {$P_2$};

    \draw[->, very thick, magenta] (0.2,0.1) -- (3.2,1.1);
    \node[above, magenta] at (1.8,0.8) {$\vec{a}$};
\end{tikzpicture}
\end{center}

\newpage

Einheitsvektor $\hat{a}$: gleiche Richtung wie $\vec{a}$

\[
\text{Betrag } |\hat{a}|=1
\]

{\Large\textcolor{red}{Multiplikation und Addition}}

1) Multiplikation mit Skalar $\to$ Streckung / Stauchung, evtl. mit Richtungsumkehr.

\[
\vec{b}=\alpha \vec{a}
\]

\[
\text{Betrag } |\vec{b}| = |\alpha|\,|\vec{a}|
\]

Richtung wie $\vec{a}$ \qquad $(\alpha > 0)$

oder entgegengesetzt \qquad $(\alpha < 0)$

Spezialfall:
\[
\vec{a}=|\vec{a}|\,\hat{a} = a\,\hat{a}
\qquad \wedge \qquad
\hat{a}=\frac{\vec{a}}{a}
\qquad (a\neq 0)
\]

Assoziativgesetz:
\[
(\alpha\beta)\vec{a}=\alpha(\beta\vec{a})
\]

\begin{center}
\begin{tikzpicture}[scale=1.1]
    % x-axis
    \draw[magenta, thick] (-0.5,0) -- (5.5,0);

    % marks 0 and 1
    \draw[magenta] (0,-0.08) -- (0,0.08);
    \node[below, magenta] at (0,-0.08) {$0$};

    \draw[magenta] (2,-0.08) -- (2,0.08);
    \node[below, magenta] at (2,-0.08) {$1$};

    % alpha label
    \node[below right, magenta] at (2.1,0) {$\alpha$};

    % green parallel lines
    \draw[green!70!black, thick] (1.6,2.1) -- (3.0,-0.8);
    \draw[green!70!black, thick] (1.2,2.3) -- (2.6,-0.6);

    % dashed ray
    \draw[red, dashed] (0,0) -- (3.6,2.9);

    % vector a
    \draw[->, very thick, red] (0,0) -- (1.4,1.25);
    \node[left, red] at (0.8,0.8) {$\vec{a}$};

    % vector b
    \draw[->, very thick, cyan] (0,0) -- (1.9,1.65);
    \node[left, cyan] at (1.3,1.45) {$\vec{b}$};

    % small green marking
    \draw[green!70!black, thick] (2.0,0.1) -- (2.2,0.9);
    \draw[green!70!black, thick] (2.15,0.05) -- (2.35,0.85);

    \node[right] at (4.0,1.3) {Strahlensatz};
\end{tikzpicture}
\end{center}

\newpage

Spezialfälle:
\[
(1)\ \alpha=0:\qquad \vec{0}=0\cdot \vec{a}
\qquad \textcolor{red}{Nullvektor}
\]

\[
(2)\ \alpha=-1:\qquad \vec{b}=-\vec{a}
\qquad \textcolor{red}{inverser Vektor}
\]

\[
(3)\ \alpha=\frac{1}{a}\ (a\neq 0):\qquad \hat{a}=\frac{\vec{a}}{a}
\qquad \textcolor{red}{Einheitsvektor}
\]

2) Addition zweier Vektoren $\to$ Hintereinanderausführung

\[
\vec{c}=\vec{a}+\vec{b}
\]

Kommutativgesetz:
\[
\vec{a}+\vec{b}=\vec{b}+\vec{a}
\]

Assoziativgesetz:
\[
\vec{a}+(\vec{b}+\vec{c})=(\vec{a}+\vec{b})+\vec{c}=\vec{a}+\vec{b}+\vec{c}
\]

Distributivgesetze:
\[
\alpha(\vec{a}+\vec{b})=\alpha\vec{a}+\alpha\vec{b}
\]

\[
(\alpha+\beta)\vec{a}=\alpha\vec{a}+\beta\vec{a}
\]

\begin{center}
\begin{tikzpicture}[scale=1.25]
    % lower vector b
    \draw[->, very thick, magenta] (0,0) -- (1.7,0);
    \node[below left, magenta] at (0.2,0) {$\vec{b}$};

    % upper vector b
    \draw[->, very thick, magenta] (1.1,1.0) -- (2.8,1.0);
    \node[above, magenta] at (2.1,1.0) {$\vec{b}$};

    % vector a
    \draw[->, very thick, magenta] (0,0) -- (1.1,1.0);
    \node[left, magenta] at (0.55,0.6) {$\vec{a}$};

    % translated a
    \draw[->, very thick, magenta] (1.7,0) -- (2.8,1.0);
    \node[right, magenta] at (2.25,0.55) {$\vec{a}$};

    % resultant c
    \draw[->, very thick, magenta] (0,0) -- (2.8,1.0);
    \node[below right, magenta] at (1.7,0.55) {$\vec{c}$};
\end{tikzpicture}
\end{center}

\newpage

{\Large\textcolor{red}{1.1.2 Skalarprodukt (inneres Produkt)}}

Def.:
\[
\vec{a}\cdot\vec{b}=ab\cos\varphi
\]

Eigenschaften:
\[
\vec{a}\cdot\vec{b}=\vec{b}\cdot\vec{a}
\qquad \text{(kommutativ)}
\]

\[
\vec{a}\cdot(\vec{b}+\vec{c})=\vec{a}\cdot\vec{b}+\vec{a}\cdot\vec{c}
\qquad \text{(distributiv)}
\]

\[
\alpha(\vec{a}\cdot\vec{b})=(\alpha\vec{a})\cdot\vec{b}
=\vec{a}\cdot(\alpha\vec{b})=\alpha\,\vec{a}\cdot\vec{b}
\qquad \text{(homogen)}
\]

Spezialfälle:
\[
\vec{a}\perp\vec{b}\ \Rightarrow\ \vec{a}\cdot\vec{b}=0
\]

\[
\vec{a}\parallel\vec{b}\ \Rightarrow\ \vec{a}\cdot\vec{b}=ab
\]

\[
\vec{a}\ \text{antiparallel zu}\ \vec{b}\ \Rightarrow\ \vec{a}\cdot\vec{b}=-ab
\]

\[
\vec{a}=\vec{b}\ \Rightarrow\ \vec{a}\cdot\vec{b}=\vec{a}^{\,2}=a^2
\]

\[
\vec{a}=\vec{0}\ \text{oder}\ \vec{b}=\vec{0}\ \Rightarrow\ \vec{a}\cdot\vec{b}=0
\]

\[
\textcolor{red}{\text{Schwarz'sche Ungleichung:}}
\qquad -ab \leq \vec{a}\cdot\vec{b} \leq ab
\]

\begin{center}
\begin{tikzpicture}[scale=1.3]
    \draw[->, very thick, magenta] (0,0) -- (2.4,0);
    \node[below right, magenta] at (2.2,0) {$\vec{b}$};

    \draw[->, very thick, magenta] (0,0) -- (1.3,1.8);
    \node[above, magenta] at (1.3,1.8) {$\vec{a}$};

    \draw[magenta] (0.7,0) arc[start angle=0,end angle=54,radius=0.7];
    \node[magenta] at (0.85,0.35) {$\varphi$};
\end{tikzpicture}
\end{center}

\end{document}